\documentclass[11pt,a4paper]{article}
\usepackage[T1]{fontenc}
\usepackage{lmodern}
\usepackage[margin=27mm]{geometry}
\usepackage{amsmath,amssymb,amsthm,mathtools}
\usepackage{microtype}
\usepackage{booktabs,array}
\usepackage{enumitem}
\usepackage[hidelinks]{hyperref}
\setlength{\parskip}{5pt plus 1pt}
\setlength{\parindent}{0pt}
\setlength{\emergencystretch}{3em}
\setlist{itemsep=2pt,topsep=3pt}
\raggedbottom
\newtheorem{theorem}{Theorem}[section]
\newtheorem{lemma}[theorem]{Lemma}
\newtheorem{corollary}[theorem]{Corollary}
\newtheorem{proposition}[theorem]{Proposition}
\newtheorem{conjecture}[theorem]{Conjecture}
\theoremstyle{definition}
\newtheorem{definition}[theorem]{Definition}
\newtheorem{example}[theorem]{Example}
\theoremstyle{remark}
\newtheorem{remark}[theorem]{Remark}
\newcommand{\R}{\mathbb{R}}
\newcommand{\N}{\mathbb{N}}
\newcommand{\Z}{\mathbb{Z}}
\newcommand{\F}{\mathbb{F}}
\newcommand{\per}{\operatorname{per}}
\newcommand{\T}{\mathsf{T}}
\newcommand{\para}[1]{\par\medskip\noindent\textbf{#1}\enspace}
\newcommand{\lean}[1]{\texttt{#1}}
\hypersetup{pdftitle={Gyires's permanent inequality for nonnegative circulant matrices of order five},pdfauthor={Dev D. Goyal},pdfsubject={Permanents of doubly stochastic matrices}}
\title{Gyires's permanent inequality for nonnegative circulant matrices of order five}
\author{Dev D. Goyal}
\date{October 2026}
\begin{document}
\maketitle
\begin{abstract}
B.~Gyires conjectured in 1978 that every doubly stochastic $n\times n$ matrix $A$ satisfies
$4\per(A)^2\ge\frac{n!}{n^n}\bigl(\per(AA^\T)+\per(A^\T A)+2\per(A^2)\bigr)$, with equality only at the matrix
$J_n$ all of whose entries are $1/n$; the conjecture is open for $n\ge5$. We prove the conjecture for every doubly
stochastic circulant matrix of order five, that is, for every convex combination of the five powers of the cyclic
permutation matrix, including the matrices with zero entries. The proof establishes the stronger inequality
$24\per(CC^\T)\le625\per(C)^2$ for every circulant $C$ of order five with nonnegative entries, with equality only when
all entries of $C$ are equal, together with a quantitative form: if $z_1,\dots,z_4$ are the gaps between consecutive
values when the entries of the first row of $C$ are sorted in increasing order, then
$625\per(C)^2-24\per(CC^\T)\ge484\,(z_1^{10}+z_2^{10}+z_3^{10}+z_4^{10})$, and the constant $484$ is best possible.
After an affine change of indices modulo five, the entries of the first row may be assumed to increase along one of six
canonical orderings, and for each of them $625\per(C)^2-24\per(CC^\T)$, written in the smallest entry and the gaps, is
a polynomial with nonnegative integer coefficients. The results have been checked by computer and proved formally in
the Lean proof assistant.
\end{abstract}
\begin{center}\begin{minipage}{0.88\linewidth}\small
\textbf{How this paper was produced.} This paper reports the result of an AI research effort directed by
Dev~D.~Goyal. ChatGPT (OpenAI), using the GPT-6 model with Pro reasoning and working autonomously from prompts by
Dev~D.~Goyal, proposed the problem, found the proof, wrote its Lean~4 formalization, and wrote a first draft of the paper.
Dev~D.~Goyal chose which problems to pursue and asked for proofs, formal verification and revisions, but did not supply
proof ideas. Claude Opus~5.5, used with Claude Code, then checked the work independently: it ported the formal proofs to
a current version of Lean, rebuilt them from source and confirmed that they rely only on Lean's three standard axioms,
tested every claim with separately written exact computations, and read every step of the proofs. It also wrote the
present version of this paper. Dev~D.~Goyal has reviewed the proofs and is responsible for them; this review is not an
independent expert validation. The paper has not been peer reviewed. The section ``About this paper'' at the end gives
the details.
\end{minipage}\end{center}

\section{Introduction}

\subsection{The conjecture}\label{sec:conjecture}
A real square matrix is \emph{doubly stochastic} if its entries are nonnegative and all its row sums and column sums
equal $1$. We write $\Omega_n$ for the set of doubly stochastic $n\times n$ matrices, $J_n\in\Omega_n$ for the matrix
all of whose entries equal $1/n$, and
\[
\per A=\sum_{\sigma\in S_n}\prod_{i=1}^n a_{i\sigma(i)}
\]
for the permanent of $A=(a_{ij})$. By van der Waerden's conjecture, proved independently by Gyires~\cite{gyires80},
Egorychev~\cite{egorychev} and Falikman~\cite{falikman}, $\per A\ge n!/n^n$ for every $A\in\Omega_n$, with equality
only at $J_n$. B.~Gyires~\cite{gyires} proposed a related lower bound for the square of the permanent. It is
Conjecture~20 in the survey of open problems on permanents by Cheon and Wanless~\cite{cw}, in the following form, where
$A^*$ denotes the conjugate transpose.

\begin{conjecture}[Gyires; Cheon and Wanless~{\cite[Conjecture~20]{cw}}]\label{conj:gyires}
For every $A\in\Omega_n$,
\[
\frac{4\per(A)^2}{\per(AA^*)+\per(A^*A)+2\per(A^2)}\ \ge\ \frac{n!}{n^n},
\]
with equality if and only if $A=J_n$.
\end{conjecture}

For real matrices $A^*=A^\T$, the transpose, and this is the form used below. For $A\in\Omega_n$ the matrices
$AA^\T$, $A^\T A$ and $A^2$ are doubly stochastic, so by van der Waerden's theorem each of their permanents is at
least $n!/n^n$, and the denominator is positive. The conjecture is therefore equivalent to the polynomial inequality
\begin{equation}\label{eq:gyires-poly}
\frac{n!}{n^n}\bigl(\per(AA^\T)+\per(A^\T A)+2\per(A^2)\bigr)\ \le\ 4\per(A)^2\qquad(A\in\Omega_n),
\end{equation}
with equality only at $J_n$. We call~\eqref{eq:gyires-poly} the \emph{homogeneous form} of Gyires's inequality: both
sides are homogeneous of degree $2n$ in the entries of $A$, so it makes sense for every square matrix.

\subsection{What this paper proves}
A matrix $C$ of order five is \emph{circulant} if its entries depend only on the difference of the indices modulo five:
$c_{ij}=w_{j-i}$ for a vector $w=(w_0,\dots,w_4)\in\R^5$, the indices being read in $\Z/5\Z$. We call $w$, the first
row of $C$, its vector of \emph{weights}, and write $C=C(w)$. A circulant is doubly stochastic exactly when $w\ge0$ and
$w_0+\dots+w_4=1$, so the doubly stochastic circulants of order five are the convex combinations $\sum_k w_kQ^k$ of
the powers of the cyclic permutation matrix $Q=C(0,1,0,0,0)$. They form a four-dimensional simplex inside the
sixteen-dimensional polytope $\Omega_5$; this simplex contains matrices with zero entries, among them its vertices
$Q^k$. For $n=5$ the constant of the conjecture is $5!/5^5=24/625$.

\para{Main results.} \emph{Let $C=C(w)$ be a circulant of order five with nonnegative weights.}
\begin{enumerate}[label=(\alph*)]
\item (Theorem~\ref{thm:strong}, the strong inequality.) $24\per(CC^\T)\le625\per(C)^2$, with equality if and only if
all five weights are equal.
\item (Theorem~\ref{thm:margin}, the margin.) If $w_{\pi(0)}\le w_{\pi(1)}\le\dots\le w_{\pi(4)}$ for a permutation
$\pi$ and $z_j=w_{\pi(j)}-w_{\pi(j-1)}$ for $1\le j\le4$, then
$625\per(C)^2-24\per(CC^\T)\ge484\,(z_1^{10}+z_2^{10}+z_3^{10}+z_4^{10})$, and no constant larger than $484$ is
possible.
\item (Corollary~\ref{cor:gyires}.) The homogeneous form~\eqref{eq:gyires-poly} of Gyires's inequality holds for $C$,
with equality if and only if all five weights are equal. In particular Conjecture~\ref{conj:gyires} holds for every
doubly stochastic circulant of order five, including its equality case.
\end{enumerate}
The strong inequality~(a) implies Gyires's inequality for $C$: for every real circulant $C$ one has $CC^\T=C^\T C$
and $\per(C^2)\le\per(CC^\T)$ (Corollary~\ref{cor:gram}), so the left-hand side of~\eqref{eq:gyires-poly} is at most
$\frac{24}{625}\cdot4\per(CC^\T)$, which by~(a) is at most $4\per(C)^2$.

\subsection{What was known}
For $n=2$ every matrix in $\Omega_2$ has the form $\bigl(\begin{smallmatrix}a&1-a\\1-a&a\end{smallmatrix}\bigr)$
with $0\le a\le1$. With $s=a^2+(1-a)^2$ and $t=2a(1-a)$ one finds $\per(A)=s$ and
$\per(AA^\T)=\per(A^\T A)=\per(A^2)=s^2+t^2$, so the conjecture reads $2s^2\ge s^2+t^2$, that is, $s\ge t$, or
$(2a-1)^2\ge0$, with equality only at $a=\frac12$, that is, at $J_2$. The inequality for $n=3$ was proved by
D.~K.~Chang~\cite{chang}. In 2026 C.~Ozols~\cite{ozols} proved the case $n=4$ for all of $\Omega_4$, including the
equality case. That proof is computer-assisted and formalized in Lean; its finite certificate is checked by compiled
code (\texttt{native\_decide}) rather than by Lean's kernel alone. Ozols's formalization also proves the inequality
$3\per(AA^\T)\le32\per(A)^2$ on $\Omega_4$, the analogue of~(a) for $n=4$, and the comparison of
Lemma~\ref{lem:gram} below. Ozols's account of prior work~\cite{ozols} also reports that Cheon and
Wanless~\cite{cw} record a sufficient condition of Chang for every $n$: the inequality holds when every entry of $A$ is
at least $(n-2)/(n-1)^2$. For $n=5$ this condition covers the doubly stochastic circulants all of whose weights lie in
$[\frac3{16},\frac14]$, a neighbourhood of $J_5$, but no circulant with a zero weight. The same account reports that a
catalogue of Minc's conjectures maintained by Y.~Lavi lists the conjecture as open and announces results, stated
without proofs, for matrices in $\Omega_4$ with a zero entry, for matrices near $J_n$ and for special families. That
catalogue could not be consulted for this paper, and whether those special families include circulants is not known. No earlier treatment of the circulant case of order five, or of the inequalities (a) and~(b), has been
found; the searches for earlier work are described in the section ``About this paper''.

\subsection{Motivation}
Gyires's conjecture is open from order five on. Circulant matrices form a natural family on which to test it, and
their permanents have long been studied (see, for instance, Minc~\cite{minc72} for doubly stochastic circulants of
the form $\alpha I+\beta Q+\gamma Q^2$): the
doubly stochastic circulants are exactly the doubly stochastic matrices invariant under the simultaneous cyclic shift
of rows and columns, every real circulant $C$ is normal, $CC^\T=C^\T C$, and the permanents in the conjecture become
polynomials in five weights. The family is four-dimensional, so the result cannot be obtained by checking finitely many
matrices, and it reaches the boundary of $\Omega_5$, so it is not covered by results for matrices near $J_5$. The
result shows that the conjecture holds on this family, with the margin~(b), which vanishes only at the constant
matrices, and the proof exhibits coordinates in which the relevant polynomial has nonnegative coefficients.

\subsection{How to read this paper}
Section~\ref{sec:background} collects the identities for circulants and a comparison lemma for permanents. Section~\ref{sec:results}
states the results, and Section~\ref{sec:proofs} proves them. Section~\ref{sec:checks} describes the computer checks
and the formal proof, and Section~\ref{sec:discussion} discusses the scope of the results.

\section{Background}\label{sec:background}

\subsection{Notation}
Indices of rows, columns and weights are taken in $\Z/5\Z=\{0,1,2,3,4\}$. For $w\in\R^5$ we write $C(w)$ for the
circulant with entries $C(w)_{ij}=w_{j-i}$, and
\[
P(w)=\per C(w),\qquad g(w)_k=\sum_{i=0}^4w_iw_{i+k},\qquad h(w)_k=\sum_{i=0}^4w_iw_{k-i}\qquad(k\in\Z/5\Z),
\]
so that $g(w)$ is the circular autocorrelation of $w$ and $h(w)$ its circular self-convolution. Note that
$g(w)_{-k}=g(w)_k$. The \emph{defects}, $D_s$ for the strong inequality and $D_G$ for Gyires's inequality, are the
homogeneous polynomials of degree ten
\begin{equation}\label{eq:defects}
D_s(w)=625\,P(w)^2-24\,P(g(w)),\qquad D_G(w)=625\,P(w)^2-12\bigl(P(g(w))+P(h(w))\bigr).
\end{equation}

\begin{lemma}\label{lem:products}
For every $w\in\R^5$, $C(w)C(w)^\T=C(w)^\T C(w)=C(g(w))$ and $C(w)^2=C(h(w))$.
\end{lemma}

\begin{proof}
Write $C=C(w)$. Substituting $m=k-j$, respectively $m=i-k$ and $m=k-i$,
\begin{align*}
(CC^\T)_{ij}&=\sum_kw_{k-i}w_{k-j}=\sum_mw_mw_{m+j-i}=g_{j-i},\\
(C^\T C)_{ij}&=\sum_kw_{i-k}w_{j-k}=\sum_mw_mw_{m+j-i}=g_{j-i},
\end{align*}
and $(C^2)_{ij}=\sum_kw_{k-i}w_{j-k}=\sum_mw_mw_{(j-i)-m}=h_{j-i}$.
\end{proof}

\begin{lemma}\label{lem:P}
For every $w=(w_0,\dots,w_4)\in\R^5$, writing $w=(a,b,c,d,e)$,
\begin{align*}
P(w)={}&a^5+b^5+c^5+d^5+e^5+15\,abcde\\
&+5\bigl(a^3be+a^3cd+a^2b^2d+a^2bc^2+a^2ce^2+a^2d^2e+ab^3c+ab^2e^2+abd^3+ac^3e\\
&\qquad+ac^2d^2+ade^3+b^3de+b^2c^2e+b^2cd^2+bc^3d+bce^3+bd^2e^2+c^2de^2+cd^3e\bigr).
\end{align*}
\end{lemma}

\begin{proof}
By definition $P(w)=\sum_{\sigma\in S_5}\prod_iw_{\sigma(i)-i}$, a sum of $120$ monomials of degree five; collecting
them gives the $26$ terms displayed. For instance, the five translations $\sigma(i)=i+k$ give $w_k^5$, and the monomial
$abcde$ arises from the $15$ permutations $\sigma$ for which $i\mapsto\sigma(i)-i$ is a bijection of $\Z/5\Z$. The
coefficients add up to $120=5!$.
\end{proof}

The following comparison is used to pass from the stronger inequality to Gyires's inequality.

\begin{lemma}\label{lem:gram}
For real $n\times n$ matrices $M$ and $N$, $\ 2\per(MN^\T)\le\per(MM^\T)+\per(NN^\T)$.
\end{lemma}

\begin{proof}
Let $x_1,\dots,x_n$ be indeterminates and define coefficients $\mu_\alpha$ and $\nu_\alpha$, indexed by
$\alpha\in\N^n$ with $\alpha_1+\dots+\alpha_n=n$, by
\[
\prod_{i=1}^n\Bigl(\sum_{k=1}^nM_{ik}x_k\Bigr)=\sum_\alpha \mu_\alpha x^\alpha,\qquad
\prod_{i=1}^n\Bigl(\sum_{k=1}^nN_{ik}x_k\Bigr)=\sum_\alpha \nu_\alpha x^\alpha .
\]
Thus $\mu_\alpha=\sum_\varphi\prod_iM_{i\varphi(i)}$, the sum over the maps $\varphi\colon\{1,\dots,n\}\to\{1,\dots,n\}$
with occupancy $\alpha$, that is, with $|\varphi^{-1}(k)|=\alpha_k$ for every $k$; similarly for $\nu_\alpha$. Expanding,
\[
\per(MN^\T)=\sum_{\sigma\in S_n}\prod_i\sum_kM_{ik}N_{\sigma(i)k}
=\sum_{\sigma\in S_n}\sum_\varphi\prod_iM_{i\varphi(i)}\prod_jN_{j,\varphi(\sigma^{-1}(j))}.
\]
For fixed $\varphi$ with occupancy $\alpha$, the map $\sigma\mapsto\varphi\circ\sigma^{-1}$ takes every map of
occupancy $\alpha$ exactly $\alpha!=\prod_k\alpha_k!$ times, and no other map. Hence
\begin{equation}\label{eq:occupancy}
\per(MN^\T)=\sum_\alpha\alpha!\,\mu_\alpha \nu_\alpha .
\end{equation}
Applying~\eqref{eq:occupancy} also to $(M,M)$ and $(N,N)$, the claim follows from $2\mu_\alpha\nu_\alpha\le
\mu_\alpha^2+\nu_\alpha^2$ and $\alpha!>0$.
\end{proof}

\begin{corollary}\label{cor:gram}
For every $w\in\R^5$, $P(h(w))\le P(g(w))$. Consequently $D_G(w)\ge D_s(w)$.
\end{corollary}

\begin{proof}
Take $M=C(w)$ and $N=C(w)^\T$ in Lemma~\ref{lem:gram}. By Lemma~\ref{lem:products}, $MN^\T=C(h(w))$ and
$MM^\T=NN^\T=C(g(w))$, so $2P(h(w))\le2P(g(w))$. Finally $D_G-D_s=12\bigl(P(g)-P(h)\bigr)$.
\end{proof}

\section{Results}\label{sec:results}

\begin{theorem}\label{thm:strong}
Let $w\in\R^5$ with $w_i\ge0$ for all $i$, and let $C=C(w)$. Then
\[
24\per(CC^\T)\le625\per(C)^2,\qquad\text{that is,}\qquad \frac{5!}{5^5}\per(CC^\T)\le\per(C)^2 .
\]
Equality holds if and only if $w_0=w_1=w_2=w_3=w_4$.
\end{theorem}

\begin{theorem}\label{thm:margin}
Let $w\in\R^5$ with $w_i\ge0$ for all $i$, let $C=C(w)$, let $\pi$ be a permutation of $\Z/5\Z$ with
$w_{\pi(0)}\le w_{\pi(1)}\le\dots\le w_{\pi(4)}$, and put $z_j=w_{\pi(j)}-w_{\pi(j-1)}$ for $1\le j\le4$. Then
\[
D_s(w)=625\per(C)^2-24\per(CC^\T)\ \ge\ 484\,\bigl(z_1^{10}+z_2^{10}+z_3^{10}+z_4^{10}\bigr).
\]
Equality holds for $w=(1,1,0,0,0)$, for which $z_1^{10}+\dots+z_4^{10}=1$ and $D_s(w)=484$. Consequently, if a real
constant $K$ satisfies $D_s(w)\ge K\,(z_1^{10}+\dots+z_4^{10})$ for all such $w$ and $\pi$, then $K\le484$.
\end{theorem}

\begin{corollary}\label{cor:gyires}
Let $w\in\R^5$ with $w_i\ge0$ for all $i$, and let $C=C(w)$. Then
\[
24\bigl(\per(CC^\T)+\per(C^\T C)+2\per(C^2)\bigr)\le2500\per(C)^2,
\]
with equality if and only if all five weights are equal. If moreover $w_0+\dots+w_4=1$, then $C\in\Omega_5$, the
denominator below is positive, and
\[
\frac{4\per(C)^2}{\per(CC^\T)+\per(C^\T C)+2\per(C^2)}\ \ge\ \frac{24}{625}=\frac{5!}{5^5},
\]
with equality if and only if $C=J_5$. Thus Conjecture~\ref{conj:gyires} holds for every doubly stochastic circulant
matrix of order five.
\end{corollary}

\section{Proofs}\label{sec:proofs}

\para{The idea in brief.} By Corollary~\ref{cor:gram} it suffices to bound $D_s$. The defect is invariant under the
affine changes of indices $i\mapsto ai+b$ of $\Z/5\Z$, and these reduce the $120$ possible orderings of the weights to
six. In each of the six chambers, the defect written in the smallest weight and the gaps between consecutive sorted
weights is a polynomial with nonnegative integer coefficients, which proves the inequality, the margin and the equality
case at once.

\subsection{Affine symmetry}
For $a\in(\Z/5\Z)^\times$ and $b\in\Z/5\Z$ let $(T_{a,b}w)_i=w_{ai+b}$; the twenty maps $i\mapsto ai+b$ are the affine
bijections of $\Z/5\Z$.

\begin{lemma}\label{lem:affine}
For all $w\in\R^5$, $a$ and $b$: $P(T_{a,b}w)=P(w)$, $g(T_{a,b}w)_k=g(w)_{ak}$ and $h(T_{a,b}w)_k=h(w)_{ak+2b}$.
Consequently $D_s(T_{a,b}w)=D_s(w)$ and $D_G(T_{a,b}w)=D_G(w)$.
\end{lemma}

\begin{proof}
We have $C(T_{a,b}w)_{ij}=w_{a(j-i)+b}=C(w)_{ai,\,aj+b}$, so $C(T_{a,b}w)$ arises from $C(w)$ by permuting the rows by
$i\mapsto ai$ and the columns by $j\mapsto aj+b$; permanents are invariant under such permutations. Substituting
$m=ai+b$,
\begin{align*}
g(T_{a,b}w)_k&=\sum_iw_{ai+b}w_{a(i+k)+b}=\sum_mw_mw_{m+ak},\\
h(T_{a,b}w)_k&=\sum_iw_{ai+b}w_{a(k-i)+b}=\sum_mw_mw_{ak+2b-m}.
\end{align*}
Hence $g(T_{a,b}w)=T_{a,0}\,g(w)$ and $h(T_{a,b}w)=T_{a,2b}\,h(w)$, and the first part gives
$P(g(T_{a,b}w))=P(g(w))$ and $P(h(T_{a,b}w))=P(h(w))$.
\end{proof}

\subsection{Six chambers}
A \emph{canonical ordering} is one of the six sequences
\begin{equation}\label{eq:orders}
01234,\quad01243,\quad01324,\quad01342,\quad01423,\quad01432,
\end{equation}
that is, an ordering $\kappa=(\kappa(0),\dots,\kappa(4))$ of $\Z/5\Z$ with $\kappa(0)=0$ and $\kappa(1)=1$. For a
canonical ordering $\kappa$ and $z=(z_0,\dots,z_4)\in\R^5$ let $L_\kappa(z)\in\R^5$ be the vector with
\[
L_\kappa(z)_{\kappa(j)}=z_0+z_1+\dots+z_j\qquad(0\le j\le4).
\]
If $z\ge0$, the weights of $L_\kappa(z)$ are nondecreasing along $\kappa$, $z_0$ is the smallest weight and
$z_1,\dots,z_4$ are the gaps between consecutive sorted weights. The \emph{chamber} of $\kappa$ is the set
$\{L_\kappa(z):z\ge0\}$ of nonnegative weight vectors that are nondecreasing along $\kappa$.

\begin{lemma}\label{lem:chambers}
Let $\pi$ be a permutation of $\Z/5\Z$. There are $a\in(\Z/5\Z)^\times$, $b\in\Z/5\Z$ and a canonical ordering
$\kappa$ with $\pi(j)=a\kappa(j)+b$ for every $j$, and the triple $(a,b,\kappa)$ is unique. If $w\in\R^5$ satisfies
$w_{\pi(0)}\le\dots\le w_{\pi(4)}$ and $w_{\pi(0)}\ge0$, then $T_{a,b}w=L_\kappa(z)$ with $z_0=w_{\pi(0)}$ and
$z_j=w_{\pi(j)}-w_{\pi(j-1)}$ $(1\le j\le4)$, and $z\ge0$.
\end{lemma}

\begin{proof}
Put $b=\pi(0)$ and $a=\pi(1)-\pi(0)\ne0$, and $\kappa(j)=a^{-1}(\pi(j)-b)$. Then $\kappa$ is an ordering of
$\Z/5\Z$ with $\kappa(0)=0$ and $\kappa(1)=1$, hence canonical, and $\pi(j)=a\kappa(j)+b$; conversely the conditions
$\kappa(0)=0$ and $\kappa(1)=1$ force $b=\pi(0)$ and $a=\pi(1)-\pi(0)$. For the second statement,
$(T_{a,b}w)_{\kappa(j)}=w_{a\kappa(j)+b}=w_{\pi(j)}=z_0+\dots+z_j$, and $z\ge0$ because the sorted weights are
nondecreasing and the smallest is nonnegative.
\end{proof}

\subsection{The chamber certificate}

\begin{lemma}\label{lem:certificate}
For each canonical ordering $\kappa$, the polynomial
\[
R_\kappa(z)=D_s\bigl(L_\kappa(z)\bigr)-484\,\bigl(z_1^{10}+z_2^{10}+z_3^{10}+z_4^{10}\bigr)
\]
in $z=(z_0,\dots,z_4)$ has nonnegative integer coefficients. More precisely, $D_s(L_\kappa(z))$ is homogeneous of degree
ten, so it has $\binom{14}{4}=1001$ coefficients; $996$ of them are positive, and the five coefficients of $z_0^{10}$ and
$z_0^9z_j$ $(1\le j\le4)$ are zero. The coefficients of $z_1^{10},z_2^{10},z_3^{10},z_4^{10}$ are
$233\,296$, $44\,209$, $484$ and $601$, and every positive coefficient is at least $484$.
\end{lemma}

\begin{proof}
This is a finite computation with integer polynomials. By Lemma~\ref{lem:P}, $P(w)$ is an explicit integer polynomial,
and by definition so is $P(g(w))$; one substitutes $w=L_\kappa(z)$ into
$625P(w)^2-24P(g(w))$, expands, and reads off the $1001$ coefficients, for each of the six orderings~\eqref{eq:orders}.
All coefficient claims of the lemma are obtained in this way. The statement about $R_\kappa$ follows: the four pure
powers $z_j^{10}$, $j\ge1$, have coefficients at least $484$, and every other coefficient is unchanged.
\end{proof}

The computation can be organized so that it is small enough to be checked by a proof assistant, and this is how it is
done in the formal proof (Section~\ref{sec:checks}). Let $u=L_\kappa(0,z_1,\dots,z_4)$, so that
$L_\kappa(z)=z_0\mathbf 1+u$, where $\mathbf 1=(1,\dots,1)$. Two identities make the dependence on $z_0$
explicit. First, $P(t\mathbf1+u)=\sum_{k=0}^5p_k(u)\,t^k$ for explicit polynomials $p_k$. Second,
$g(t\mathbf1+u)=(5t^2+2tS)\mathbf1+g(u)$, where $S=u_0+\dots+u_4$, so that
$P(g(t\mathbf1+u))=\sum_{k=0}^5p_k(g(u))\,(5t^2+2tS)^k$. Hence $D_s(L_\kappa(z))=\sum_{j=0}^{10}Q_j\,z_0^{\,j}$ with
explicit polynomials $Q_j$ in $z_1,\dots,z_4$, and each $Q_j$ is expanded in powers of $z_1$ with coefficients that are
polynomials in $z_2,z_3,z_4$; by Lemma~\ref{lem:certificate}, $Q_9=Q_{10}=0$. The nonnegativity of the coefficients of
$R_\kappa$ amounts to the statement that, after $484\,(z_1^{10}+\dots+z_4^{10})$ is subtracted from $Q_0$, each of
these polynomials in $z_2,z_3,z_4$ has nonnegative coefficients.

\subsection{Proof of Theorems~\ref{thm:margin} and~\ref{thm:strong}}

\begin{proof}[Proof of Theorem~\ref{thm:margin}]
Let $w\ge0$ and $\pi$ be as in the theorem, and choose $a$, $b$, $\kappa$ and $z$ by Lemma~\ref{lem:chambers}. By
Lemmas~\ref{lem:affine} and~\ref{lem:certificate}, since $z\ge0$,
\[
D_s(w)=D_s(T_{a,b}w)=D_s\bigl(L_\kappa(z)\bigr)=484\,\bigl(z_1^{10}+\dots+z_4^{10}\bigr)+R_\kappa(z)
\ \ge\ 484\,\bigl(z_1^{10}+\dots+z_4^{10}\bigr).
\]
For the equality case, let $w=(1,1,0,0,0)$, sorted by $\pi=(2,3,4,0,1)$, with gaps $(z_1,z_2,z_3,z_4)=(0,0,1,0)$. Here
$g(w)=(2,1,0,0,1)$, and Lemma~\ref{lem:P} gives $P(w)=2$ and $P(g(w))=84$, so $D_s(w)=625\cdot4-24\cdot84=484$. A
constant $K$ as in the theorem therefore satisfies $K\cdot1\le484$.
\end{proof}

\begin{proof}[Proof of Theorem~\ref{thm:strong}]
By Lemma~\ref{lem:products}, $\per(CC^\T)=P(g(w))$, so the inequality is $D_s(w)\ge0$, which follows from
Theorem~\ref{thm:margin} applied with any permutation $\pi$ that sorts $w$. If $D_s(w)=0$, Theorem~\ref{thm:margin}
gives $z_1=\dots=z_4=0$, so all weights are equal. Conversely, if every weight equals $c$, then $P(w)=120c^5$ (each of
the $120$ permutations contributes $c^5$) and $g(w)=(5c^2,\dots,5c^2)$, so $P(g(w))=120\cdot5^5c^{10}$ and
$625P(w)^2=9\,000\,000\,c^{10}=24P(g(w))$.
\end{proof}

\subsection{Proof of Corollary~\ref{cor:gyires}}

\begin{proof}
By Lemma~\ref{lem:products},
$\per(CC^\T)+\per(C^\T C)+2\per(C^2)=2P(g(w))+2P(h(w))$, so the first inequality reads $4D_G(w)\ge0$. By
Corollary~\ref{cor:gram} and Theorem~\ref{thm:strong}, $D_G(w)\ge D_s(w)\ge0$. If $D_G(w)=0$, then $D_s(w)=0$, and
all weights are equal by Theorem~\ref{thm:strong}; conversely, for constant $w$ one has $h(w)=g(w)$, hence
$D_G(w)=D_s(w)=0$.

Now let also $\sum_iw_i=1$. Every row and every column of $C$ contains each weight once, so $C\in\Omega_5$. Some weight
is positive, so $g(w)_0=\sum_iw_i^2>0$. The matrix $C(g(w))$ has nonnegative entries and diagonal entries $g(w)_0$, so
$P(g(w))\ge g(w)_0^5>0$, the identity permutation contributing $g(w)_0^5$ and the others nonnegative terms. Since also
$P(h(w))\ge0$, the denominator $2P(g(w))+2P(h(w))$ is positive, and dividing the first inequality by $625$ times the denominator gives
the ratio inequality. Equality holds if and only if all weights are equal, that is, equal to $1/5$, which means
$C=J_5$.
\end{proof}

\section{Checking the results}\label{sec:checks}

A mathematical proof can contain mistakes, and this one rests on a computation too large to be checked by hand, so the
results were checked in two independent ways.

\para{Exact computations.}
A script for the PARI/GP computer algebra system performs $72$ checks in exact integer and rational arithmetic. It was
written separately from the Lean files and does not use them, or the program that generated the certificate files; it is the file \texttt{independent\_checks.gp} in the directory
\texttt{verification} of the accompanying repository~\cite{repo}. It computes every permanent from the definition, as a
sum over all permutations, and compares it with PARI's implementation of Ryser's formula. It verifies
Lemmas~\ref{lem:products}, \ref{lem:P} and~\ref{lem:affine} symbolically, and that each of the $120$ orderings is the image of exactly one canonical ordering
under exactly one affine map (Lemma~\ref{lem:chambers}). It recomputes the chamber polynomials of
Lemma~\ref{lem:certificate}, not only for the six canonical orderings but for all $120$ orderings, which gives
$120\cdot1001=120\,120$ coefficients, of which $119\,520$ are positive and $600$ zero, with the zero coefficients, the
least positive coefficient and the four pure-power coefficients as stated. It tests Lemma~\ref{lem:gram} on $200$ random
rational matrices of orders one to five and identity~\eqref{eq:occupancy} on $24$ random integer pairs; it tests
Theorems~\ref{thm:strong} and~\ref{thm:margin} and Corollary~\ref{cor:gyires}, with their equality cases, on
$2000$ nonnegative weight vectors, many of them with zero or repeated weights; and it checks the example of
Theorem~\ref{thm:margin}, the case $n=2$ of the introduction and the remarks of Section~\ref{sec:discussion}. The
random tests are evidence, not proofs; the chamber computation, by contrast, is exact and exhaustive.

\para{Formal proof in Lean.}
Lean~\cite{lean} is a \emph{proof assistant}: a program in which mathematical statements and their proofs are written
in a precise formal language, and every step is checked mechanically down to the basic rules of logic. The results of
Section~\ref{sec:results} are proved in Lean~4 (version 4.35.0) on top of Mathlib~\cite{mathlib}, Lean's library of
formalized mathematics. The library is in the directory \texttt{lean/Circulant5} of the repository~\cite{repo}, in the
namespace \texttt{Circulant5}. There \texttt{circ w} is the circulant with entries \texttt{w (j - i)}, \texttt{P5} the
polynomial of Lemma~\ref{lem:P}, and \texttt{gram}, \texttt{conv} the vectors $g$ and $h$; \texttt{Fin 5}, with its
arithmetic modulo five, plays the role of $\Z/5\Z$.

The principal statements are restated, using only Mathlib's definitions, in a \emph{Challenge} file, which states the
theorems with their proofs omitted, the file
\texttt{lean/\allowbreak{}Challenges/\allowbreak{}GyiresCirculant5.lean} (namespace \texttt{GyiresCirculant5}). There a circulant is a matrix of
the form \texttt{Matrix.circulant v}, whose entries are \texttt{v (i - j)}; this is the transpose of $C(v)$, which
describes the same set of matrices. Doubly stochastic matrices are Mathlib's \texttt{doublyStochastic}, the permanent is
Mathlib's \texttt{Matrix.permanent}, $5!/5^5$ is the real quotient \texttt{(Nat.factorial 5 :\ $\R$) / 5 \textasciicircum{} 5},
and $J_5$ is the constant matrix with entries $5^{-1}$. The \emph{Solution} file
\texttt{lean/\allowbreak{}Solutions/\allowbreak{}GyiresCirculant5.lean} proves the same theorems from the library. Lean's
Comparator tool confirms that the theorems proved there have exactly the statements of the Challenge file and use only
the standard axioms, and the proofs are accepted by Lean's kernel and by two independent kernels, NanoDa and
con-ron. The table lists where each statement of the paper is
proved; names in the first group are in the namespace \texttt{GyiresCirculant5}, the others in \texttt{Circulant5}.
\begin{center}\small\renewcommand{\arraystretch}{1.25}
\begin{tabular}{>{\raggedright\arraybackslash}p{0.30\linewidth}>{\raggedright\arraybackslash}p{0.64\linewidth}}
\toprule
Statement in this paper & Lean declaration(s)\\
\midrule
Theorem~\ref{thm:strong} & \lean{strong\_inequality}, \lean{strong\_inequality\_eq\_iff} \\
Theorem~\ref{thm:margin} & \lean{strong\_margin}, \lean{strong\_margin\_attained}, \lean{strong\_margin\_sharp} \\
Corollary~\ref{cor:gyires}, homogeneous form & \lean{gyires\_homogeneous}, \lean{gyires\_homogeneous\_eq\_iff} \\
Corollary~\ref{cor:gyires}, ratio form & \lean{gyires\_ratio}, \lean{gyires\_ratio\_eq\_iff}, \lean{gyires\_denominator\_pos} \\
\midrule
Lemma~\ref{lem:products} & \lean{circ\_mul\_transpose}, \lean{transpose\_mul\_circ}, \lean{circ\_mul\_self} \\
Lemma~\ref{lem:P} & \lean{permanent\_circ} \\
Lemma~\ref{lem:gram}, Corollary~\ref{cor:gram} & \lean{PermanentGram.permanent\_gram\_add\_ge}, \lean{conv\_permanent\_le\_gram} \\
Lemma~\ref{lem:affine} & \lean{P5\_affine}, \lean{gram\_affine}, \lean{conv\_affine}, \lean{strongDefect\_affine}, \lean{gyiresDefect\_affine} \\
Lemma~\ref{lem:chambers} & \lean{order\_affine\_cover}, \lean{order\_affine\_cover\_unique}, \lean{canonical\_chamber\_of\_monotone}, \lean{exists\_canonical\_chamber} \\
Lemma~\ref{lem:certificate}, nonnegativity of $R_\kappa$ & \lean{strong\_chamber\_0}, \dots, \lean{strong\_chamber\_5}, \lean{strong\_chambers} \\
Theorem~\ref{thm:margin}, the example & \lean{strongDefect\_twoSupport}, \lean{strongDefect\_margin\_attained}, \lean{strongDefect\_margin\_sharp} \\
Corollary~\ref{cor:gyires}, positive denominator & \lean{positiveDenominator} \\
Section~\ref{sec:discussion}, the two-point example & \lean{P5\_twoSupport}, \lean{P5\_gram\_twoSupport}, \lean{P5\_conv\_twoSupport}, \lean{gyiresDefect\_twoSupport}, \lean{gyiresDenominator\_twoSupport},  \lean{permanent\_twoSupport\_half}, \lean{permanent\_gram\_twoSupport\_half}, \lean{permanent\_square\_twoSupport\_half}, \lean{normalizedGyiresRatio\_twoSupport\_half} \\
\bottomrule
\end{tabular}
\end{center}
The Lean project was rebuilt from source, and the axioms on which each of these theorems depends were listed with
Lean's \texttt{\#print axioms} command. Every one depends only on the three standard axioms of Lean's logic
(\texttt{propext}, \texttt{Classical.choice}, \texttt{Quot.sound}); no step is left unproved, and the proofs use neither
\texttt{native\_decide} nor any other evaluation by compiled code. The formal proof of Lemma~\ref{lem:certificate}
follows the organization described after it: the six chamber polynomials are expanded in $z_0$ and $z_1$, and each of
the resulting polynomials in $z_2,z_3,z_4$ is written with explicit nonnegative integer coefficients, proved equal to the
corresponding expansion coefficient by the \texttt{ring} tactic and nonnegative by the \texttt{positivity} tactic. These explicit
coefficients were produced by a computer program, but Lean does not rely on that program: it checks every identity and
every inequality itself. The formal proof of Lemma~\ref{lem:gram} is not the occupancy argument given here but a
sum-of-squares argument over permutations, adapted from the formalization of Ozols~\cite{ozols}.

Several remarks are also formalized: the case $n=2$ of the introduction (\lean{gyires\_two}); the characterization of
the doubly stochastic circulants of order five and their description as the combinations $\sum_kw_kQ^k$
(\lean{circulant\_mem\_doublyStochastic\_iff}, \lean{circulant\_eq\_sum\_pow}); and, for real square matrices of
every order, the reduction of Gyires's inequality to the strong form described in Section~\ref{sec:discussion}
(\lean{gyires\_of\_strong}).

\para{What is not formalized.}
The formal proof establishes $R_\kappa(z)\ge0$ for all $z\ge0$, which is the part of Lemma~\ref{lem:certificate} that
the theorems need. The finer statements of that lemma (the numbers $996$ and $5$ of positive and zero coefficients,
the positions of the zeros, the four pure-power coefficients, and the least positive coefficient $484$) are not stated
in Lean; they are verified by the exact computation described above. The remaining remarks of
Section~\ref{sec:discussion}, which rest on random tests or on the literature, are not formalized.

\section{Discussion}\label{sec:discussion}

\para{Scope.} The results concern circulant matrices of order five only. They do not prove Gyires's conjecture for all of
$\Omega_5$, or for any order $n\ge6$. Within the circulants of order five they cover every matrix in the
four-dimensional simplex of doubly stochastic circulants, including all of its faces, and every nonnegative circulant
without normalization.

\para{The coordinates.} Written in the original weights, $D_s$ has negative coefficients, so its nonnegativity is not
visible without the change to sorted gaps. In the chamber coordinates all coefficients are nonnegative, and the
zero coefficients are exactly those of the monomials $z_0^{10}$ and $z_0^9z_j$, which reflects the fact that $D_s$
vanishes to second order along the constant vectors.

\para{The two-point example.} For $w=(t,t,0,0,0)$ one has $g(w)=(2t^2,t^2,0,0,t^2)$, $h(w)=(t^2,2t^2,t^2,0,0)$,
$P(w)=2t^5$, $P(g(w))=P(h(w))=84t^{10}$, and $D_s(w)=D_G(w)=484t^{10}$; this is the example showing that $484$
cannot be improved. For $t=\frac12$ the matrix is doubly stochastic, $\per C=\frac1{16}$ and $\per(CC^\T)=\per(C^2)=\frac{21}{256}$,
and Gyires's ratio equals $\frac1{21}\approx0.048$, above $\frac{24}{625}=0.0384$. The example also shows that the
inequality $\per(AA^\T)\le\per(A)$ fails for some doubly stochastic circulants. If it held, the strong form
$\frac{n!}{n^n}\per(AA^\T)\le\per(A)^2$ would follow at once from van der Waerden's bound $\per(A)\ge n!/n^n$.
According to Ryser's account, in the translation of Best and Wanless~\cite{bw}, an inequality of this form was once
conjectured and was disproved by M.~Newman with the order-four analogue $\frac12(I+Q)$ of this example, for which
$\per A=\frac18<\frac9{64}=\per(AA^\T)$.

\para{A stronger question.} For every real square matrix $A$, Lemma~\ref{lem:gram} with $M=A$ and $N=A^\T$ gives
$2\per(A^2)\le\per(AA^\T)+\per(A^\T A)$. Hence the strong form, applied to $A$ and to $A^\T$, implies Gyires's
inequality for $A$; this is the reduction used by Ozols~\cite{ozols} for $n=4$. It is therefore natural to ask whether
\[
\frac{n!}{n^n}\per(AA^\T)\le\per(A)^2\qquad\text{for every }A\in\Omega_n .
\]
This holds for $n=2$ by the computation in the introduction, for $n=4$ by Ozols's theorem, and for the circulants of
order five by Theorem~\ref{thm:strong}. In random tests on $2500$ doubly stochastic matrices of orders $2\le n\le6$,
each a convex combination of random permutation matrices, no counterexample was found.

\para{Other orders.} Every real circulant of order $n$ is normal, and Lemma~\ref{lem:gram} gives
$\per(C^2)\le\per(CC^\T)$ in every order, so Gyires's conjecture for circulants of order $n$ follows from the strong
form $\frac{n!}{n^n}\per(CC^\T)\le\per(C)^2$. Random tests on $1080$ nonnegative circulants of orders $2\le n\le8$
found no counterexample to this strong form. The method of this paper depends essentially on the value of $n$: the
affine group of $\Z/n\Z$ reduces the $n!$ orderings of the weights to $(n-2)!$ chambers when $n$ is prime, and the
degree of the defect grows with $n$.

\section*{About this paper}\label{sec:about}
\para{How the work was done.}
This paper is the result of an AI research effort directed by Dev~D.~Goyal; the repository~\cite{repo} documents the
work. The roles were as follows.
\begin{itemize}
\item \emph{Dev~D.~Goyal} started the project, chose which problems to pursue, asked for proofs, formal verification and
revisions, and decided what to publish. Dev~D.~Goyal did not supply proof strategies or intermediate arguments.
\item \emph{ChatGPT} (OpenAI), using the GPT-6 model with Pro reasoning, worked autonomously from those prompts. It proposed
the problem, found the theorems and their proofs, including the chamber certificate, wrote the Lean~4 formalization of
the library \texttt{Circulant5} and the programs that generated its certificate files, and wrote a first draft of this
paper.
\item \emph{Claude Opus~5.5}, used with Claude Code, checked the work independently. It ported the formalization from
Lean~4.19 to Lean~4.35, adapting it to Lean's module system, rebuilt it from source, and confirmed that the theorems
listed in Section~\ref{sec:checks} rely only on Lean's standard axioms. It wrote the Challenge and Solution files and the
Lean files \texttt{Circulant5/\allowbreak{}MathlibStatements.lean}, which proves the margin of
Theorem~\ref{thm:margin} for an arbitrary sorting permutation, its attainment and sharpness in that form, and the
uniqueness in Lemma~\ref{lem:chambers} (the earlier formalization stated the margin only in the six canonical
chambers), and \texttt{Circulant5/\allowbreak{}Remarks.lean}, which formalizes the remarks listed in
Section~\ref{sec:checks}. It also wrote the PARI/GP checks and read every step of the proofs. Its checks are
independent of ChatGPT's work and of the program that generated the certificate files, but not of its own Lean files
and of this paper, which it wrote. It then wrote the
present version of the paper, including the proofs of Lemmas~\ref{lem:products}, \ref{lem:affine}
and~\ref{lem:chambers} in the form given here, the computation for $n=2$ in the introduction, and the Discussion
section. The theorems of Section~\ref{sec:results} and the proof strategy are ChatGPT's.
\end{itemize}
Dev~D.~Goyal has reviewed the proofs and is responsible for this paper; this review is not an independent expert
validation.

\para{Review status.} This paper has not been peer reviewed, and it has not been submitted to a journal.

\para{Priority.} The searches for earlier work covered the works cited here, as far as they could be accessed,
general web searches for Gyires's conjecture and for permanent inequalities for circulant matrices, the account of prior
work in Ozols's repository~\cite{ozols}, and the Lean collections google-deepmind/formal-conjectures (commit
\texttt{37f0808}, 24~October~2026), openai/math (commit \texttt{1581d4c}, 12~October~2026) and plby/lean-proofs (commit
\texttt{8822f7d}, 15~September~2026); none of them contains the circulant case. Gyires's original article, the full texts of
the survey of Cheon and Wanless and of Chang's article, and Lavi's catalogue could not be accessed; the statement of
Conjecture~\ref{conj:gyires} and the results reported from them were checked against the account of prior work in
Ozols's repository~\cite{ozols}, against the publishers' abstracts, and against other citations of these works. These
searches were limited, and no claim of priority or novelty is made.

\para{Corrections.} Please report any error in the issue tracker of the repository~\cite{repo}, where the Lean sources and
the numerical checks are also available.

\section*{Acknowledgements}
The conjecture is due to B.~Gyires. The formal proof of Lemma~\ref{lem:gram} is adapted from the Lean formalization of
the case $n=4$ by C.~Ozols~\cite{ozols}, which is licensed under the Apache License~2.0. Thanks to the developers of
Lean and Mathlib.

\begin{thebibliography}{99}\small
\raggedright
\bibitem{bw} D.~Best and I.~M.~Wanless, What did Ryser conjecture?, preprint, 2018. arXiv:1801.02893.
\bibitem{chang} D.~K.~Chang, On two permanental conjectures, \emph{Linear and Multilinear Algebra} \textbf{26} (1990),
207--213.
\bibitem{cw} G.-S.~Cheon and I.~M.~Wanless, An update on Minc's survey of open problems involving permanents,
\emph{Linear Algebra and its Applications} \textbf{403} (2005), 314--342.
\bibitem{egorychev} G.~P.~Egorychev, The solution of van der Waerden's problem for permanents, \emph{Soviet
Mathematics Doklady} \textbf{23} (1981), 619--622.
\bibitem{falikman} D.~I.~Falikman, Proof of the van der Waerden conjecture regarding the permanent of a doubly
stochastic matrix, \emph{Matematicheskie Zametki} \textbf{29} (1981), 931--938.
\bibitem{gyires80} B.~Gyires, The common source of several inequalities concerning doubly stochastic matrices,
\emph{Publicationes Mathematicae Debrecen} \textbf{27} (1980), 291--304.
\bibitem{gyires} B.~Gyires, On permanent inequalities, in \emph{Combinatorics} (Proc.\ Fifth Hungarian Colloq.,
Keszthely, 1976), A.~Hajnal and V.~T.~S\'os (eds.), Colloq.\ Math.\ Soc.\ J\'anos Bolyai \textbf{18}, North-Holland,
1978, 471--484.
\bibitem{lean} L.~de~Moura and S.~Ullrich, The Lean~4 theorem prover and programming language, in \emph{Automated
Deduction -- CADE 28}, Lecture Notes in Computer Science 12699, Springer, 2021, 625--635.
\bibitem{minc72} H.~Minc, On permanents of circulants, \emph{Pacific Journal of Mathematics} \textbf{42} (1972),
477--484.
\bibitem{mathlib} The mathlib Community, The Lean mathematical library, in \emph{Proceedings of the 9th ACM SIGPLAN
International Conference on Certified Programs and Proofs (CPP 2020)}, 367--381. Version used: commit
\texttt{56a005855ea1} (v4.35.0).
\bibitem{ozols} C.~Ozols, \emph{Gyires' permanent inequality for $4\times4$ doubly stochastic matrices}, Lean
formalization and documentation, 2026. \url{https://github.com/cozuya/gyires-permanent-4x4}, commit
\texttt{244f48c}.
\bibitem{repo} D.~D.~Goyal, \emph{Gyires's permanent inequality for circulant matrices of order five}, repository, 2026.
\url{https://github.com/dvgyl/gyires-circulant-5}.
\end{thebibliography}
\end{document}
